\section{Taichi (太极): un lenguaje de programación de una sola persona}

\subsection{Qué es la computación gráfica}

Hu Yuanming (胡渊鸣) define la computación gráfica de forma concisa y contundente: \textbf{crear un mundo virtual con sumas, restas, multiplicaciones y divisiones}. La CPU y la GPU solo pueden hacer sumas, restas, multiplicaciones, divisiones y funciones trigonométricas, pero con estos building blocks más básicos hay que crear imágenes de videojuego que parezcan vivas.

El desarrollo de la computación gráfica se divide, a grandes trazos, en varias etapas:

\begin{enumerate}
\item \textbf{Visualización y renderizado con realismo} (antes de aproximadamente 1995)
\item \textbf{Simulación física} (después de 1995): la simulación de humo de Fedkiw et al. (Visual Simulation of Smoke) abrió la era de la simulación
\item \textbf{Fotografía computacional}: técnicas como super resolution, indisociables de la visión por computadora
\item \textbf{Neural Graphics}: nuevos paradigmas impulsados por IA como DLSS, NeRF y Gaussian Splatting
\end{enumerate}

\begin{knowledgebox}{Señal para saber si un campo está en declive}
"Cuando notas que un problema se vuelve cada vez más limpio y cada vez más well-defined, eso indica que ya no queda mucho por hacer con ese problema. Cuando se forma un benchmark particularmente sistemático, cuando las soluciones se vuelven cada vez más complejas, eso indica que las oportunidades de abrir nuevo territorio ya son escasas." Hu Yuanming (胡渊鸣) ilustra esto con el ejemplo de AMCMCPPM (muestreo unificado de caminos por cadenas de Markov Monte Carlo adaptativas) en el campo del renderizado.
\end{knowledgebox}

\subsection{El nacimiento y la misión de Taichi}

La misión central de Taichi es construir una \textbf{infraestructura (Infrastructure)} que haga que la simulación sea más fácil de desarrollar, se ejecute más rápido, use menos memoria y se integre bien con el ecosistema circundante (Python, PyTorch).

Un logro emblemático fue ejecutar una simulación de mil millones ($10^9$) de partículas en una sola GPU 3090 (24 GB de memoria de video). Para lograrlo, el espacio de almacenamiento que ocupa cada partícula no puede superar los 24 bytes. Él hizo optimizaciones de cuantización extremas a nivel del compilador: las coordenadas XYZ de una partícula se representan con un entero de 32 bits (X con 11 bits, Y con 11 bits, Z con 10 bits); ese es el trabajo de QuantTaichi.

\begin{importantbox}{Por qué Taichi necesitaba construir su propio compilador}
El punto clave es que, aunque CUDA permite a los programadores de C++ escribir programas de GPU, la madurez del ecosistema de Python, de la infraestructura del compilador LLVM y de métodos de simulación como MPM creó la oportunidad de construir lenguajes específicos de dominio en un nivel más alto. Taichi se basa en la sintaxis de Python y, mediante un compilador construido desde cero, compila código de alto nivel a programas de GPU de alto rendimiento, de modo que el usuario no tiene que manejar manualmente operaciones de bits ni optimización de memoria.
\end{importantbox}

La experiencia de desarrollar Taichi de forma independiente en el MIT es descrita por Hu Yuanming (胡渊鸣) como "muy placentera": él desempeñaba a la vez cuatro roles:

\begin{table}[H]
\centering
\caption{Los cuatro roles en el desarrollo de Taichi}
\begin{tabular}{ll}
\toprule
\textbf{Rol} & \textbf{Trabajo concreto} \\
\midrule
Gerente de producto & Analizar las necesidades de los usuarios de simulación (porque él mismo era el usuario) \\
Científico & Conocer la tecnología de vanguardia y convertirla en características del producto \\
Ingeniero & CI/CD, arquitectura de software, optimización de rendimiento en C++, compilación multihilo, caché de compilación \\
CEO & Integrar los resultados de los otros tres roles para que más gente lo usara (promoción en Zhihu, etc.) \\
\bottomrule
\end{tabular}
\end{table}

"Usar mi propio lenguaje de programación para implementar mi propio programa y descubrir que además era más cómodo de usar que CUDA: creo que en ese momento me sentí un poco increíble." Este fue el momento que más lo emocionó durante el desarrollo de Taichi.

\subsection{Frozen en 99 líneas de código}

El demo más conocido de Taichi es el efecto de "Frozen en 99 líneas de código": con un código de Python extremadamente corto se implementó una simulación de nieve impulsada por el método de puntos materiales (MPM). Este demo obtuvo una atención enorme en redes sociales y se convirtió en una de las etiquetas más emblemáticas de Hu Yuanming (胡渊鸣).

Sin embargo, su actitud ante esta etiqueta es ambivalente: "Creo que hace tiempo dejaste de idealizar este tipo de etiquetas, y una etiqueta así de ninguna manera muestra quién eres en realidad." De hecho, el demo de 99 líneas es solo la punta del iceberg de lo que Taichi puede hacer: detrás hay toda una pila de compilador, un sistema de tipos, un motor de diferenciación automática, generadores de código para múltiples backends y mucha más infraestructura.

\begin{knowledgebox}{La programación diferenciable (Differentiable Programming) de Taichi}
Una característica clave de Taichi es que soporta la programación diferenciable, es decir, tu programa de simulación puede derivarse automáticamente. Esto significa que puedes:
\begin{itemize}
\item Dado un estado objetivo, optimizar de forma inversa los parámetros de la simulación
\item Incrustar el simulator en el ciclo de entrenamiento de una neural network
\item Usar gradient descent para optimizar la estrategia de control de un robot
\end{itemize}
Los dos artículos DiffTaichi y ChainQueen se basan justamente en esta característica para combinar simulation y robotics. Esta dirección fue sugerida por Jiajun (家俊), un compañero mayor de Tsinghua: "¿Por qué no intentas hacer que tu simulator sea differentiable?"
\end{knowledgebox}

Hu Yuanming (胡渊鸣) recuerda que basta con agregar una o dos líneas al código de Taichi para convertir una simulación ordinaria en una simulación diferenciable: "También pienso que solo construyendo tu propio compilador puedes lograr esto." Esta capacidad de abstracción a nivel de compilador es algo a lo que el código de CUDA escrito a mano no puede llegar.

\subsection{La visión de Taichi 2.0}

La entrevista reveló que la planeación de Taichi 2.0 está en marcha. El objetivo es que Taichi siga teniendo un uso amplio en la era de los modelos grandes, con una influencia que supere en más de 50 veces la de la versión 1.0. Las direcciones concretas podrían incluir:

\begin{itemize}
\item \textbf{Soporte para la optimización de Transformer}: si en su momento Taichi se hubiera orientado antes hacia esta dirección, "quizás ahora Taichi estaría en la posición de Triton"
\item \textbf{Hybrid Simulator}: combinar de verdad simulation y neural network, con una parte de simulación física y una parte data-driven
\item \textbf{Soporte para el ecosistema de LLM}: encontrar el escenario de aplicación de Taichi en el ámbito de los LLM
\end{itemize}

Un dato interesante: el número de GitHub Stars de Taichi supera al de proyectos como NVIDIA Warp, Triton y Mojo: "Taichi sigue siendo bastante popular, solo que no ha encontrado un ámbito grande donde desplegarse."

\subsection{Las tres evoluciones de la filosofía del código abierto}

La visión de Hu Yuanming (胡渊鸣) sobre el código abierto pasó por tres etapas:

\textbf{Primera etapa (durante el doctorado en el MIT)}: el código abierto es lo mejor. En ese entonces muchos artículos de Graphics no publicaban código abierto, y él insistió en que todos los resultados experimentales de Taichi pudieran reproducirse con un solo comando.

\textbf{Segunda etapa (inicio del emprendimiento)}: comercializar el código abierto es sumamente arriesgado. "Ya le diste a la gente lo mejor que tenías, entonces, ¿cómo vas a ganar dinero?" La inversión en Taichi disminuyó y el equipo se trasladó a trabajar en Meshy.

\textbf{Tercera etapa (actual)}: el código abierto es una herramienta estratégica de marketing.

\begin{importantbox}{El código abierto como mercadotecnia técnica}
"Es cierto que el código abierto no genera ganancias, pero un buen proyecto de código abierto puede ayudarte a atraer al mejor talento y a los mejores clientes: se convierte en un medio de marketing." En la industria de los videojuegos, la inversión en investigación y desarrollo y en promoción de mercado suele estar en una proporción de 1:1 o incluso mayor. Hacer mercadotecnia técnica a través del código abierto es el mejor camino para los fundadores técnicos que no son hábiles en el marketing tradicional. Pero esto requiere ser "muy strategic": hay que pensar en una buena forma de comercialización que lo complemente.
\end{importantbox}

Él reconoce con franqueza que el código abierto es una parte de su ideal como Founder de la que no puede desprenderse: "tarde o temprano es posible que vuelvas a tomar este camino del código abierto." Actualmente el equipo está planeando Taichi 2.0, posicionado como "infrastructure supporting Graphics and AI".

\subsection{Resumen de la sección}

El proyecto Taichi encarna el espíritu geek de Hu Yuanming (胡渊鸣) de "construir las cosas desde la capa más básica" y también es el puente que lo llevó de la academia al emprendimiento. En lo técnico, mostró el enorme potencial de un compilador específico de dominio (simulación de mil millones de partículas); en lo metodológico, forjó una capacidad integral de "gerente de producto + científico + ingeniero + CEO". Las tres evoluciones de la filosofía del código abierto reflejan, a su vez, el equilibrio entre el ideal y la realidad comercial.
